\documentclass[10pt,a4paper,twocolumn]{article}

% ============================================
% PACKAGES
% ============================================

% --- Typography and Layout ---
\usepackage[utf8]{inputenc}
\usepackage[T1]{fontenc}
\usepackage{microtype}
\usepackage{times}  % Times font (standard for ML papers)

% --- Page Layout (NeurIPS/arXiv style) ---
\usepackage[
    top=1in,
    bottom=1in,
    left=0.75in,
    right=0.75in,
    columnsep=0.25in
]{geometry}

% --- Text Justification ---
\usepackage{ragged2e}
\justifying
\setlength{\parindent}{1em}
\setlength{\parskip}{0pt}

% --- Mathematics ---
\usepackage{amsmath}
\usepackage{amssymb}
\usepackage{amsthm}
\usepackage{mathtools}
\usepackage{bm}  % Bold math symbols

% --- Graphics and Figures ---
\usepackage{graphicx}
\usepackage{subcaption}
\usepackage{float}

% --- Tables ---
\usepackage{booktabs}
\usepackage{multirow}
\usepackage{tabularx}

% --- Colors and Hyperlinks ---
\usepackage[dvipsnames]{xcolor}
\usepackage[colorlinks=true,linkcolor=MidnightBlue,citecolor=ForestGreen,urlcolor=MidnightBlue]{hyperref}

% --- Citations ---
\usepackage[numbers,sort&compress]{natbib}

% --- Algorithms ---
\usepackage{algorithm}
\usepackage{algorithmic}

% --- Miscellaneous ---
\usepackage{enumitem}
\usepackage{cleveref}

% ============================================
% CUSTOM COMMANDS
% ============================================

% Math operators
\DeclareMathOperator*{\argmax}{arg\,max}
\DeclareMathOperator*{\argmin}{arg\,min}
\DeclareMathOperator{\E}{\mathbb{E}}
\DeclareMathOperator{\Var}{Var}
\DeclareMathOperator{\KL}{KL}

% Common symbols
\newcommand{\R}{\mathbb{R}}
\newcommand{\N}{\mathbb{N}}
\newcommand{\Z}{\mathbb{Z}}
\newcommand{\x}{\mathbf{x}}
\newcommand{\y}{\mathbf{y}}
\newcommand{\z}{\mathbf{z}}
\newcommand{\w}{\mathbf{w}}
\newcommand{\h}{\mathbf{h}}
\newcommand{\thetab}{\boldsymbol{\theta}}
\newcommand{\mub}{\boldsymbol{\mu}}
\newcommand{\sigmab}{\boldsymbol{\sigma}}
\newcommand{\epsilonb}{\boldsymbol{\epsilon}}

% Theorem environments
\newtheorem{theorem}{Theorem}
\newtheorem{lemma}[theorem]{Lemma}
\newtheorem{proposition}[theorem]{Proposition}
\newtheorem{corollary}[theorem]{Corollary}
\newtheorem{definition}{Definition}
\newtheorem{remark}{Remark}

% Custom abstract with centered title
\renewenvironment{abstract}{%
    \begin{center}
        \textbf{Abstract}
    \end{center}
    \vspace{-0.5em}
    \noindent\ignorespaces
}{\par\vspace{1em}}

% ============================================
% TITLE AND AUTHORS
% ============================================

\title{\textbf{Your Paper Title Here: \\ A Modern Template for Machine Learning Research}}

\author{
    First Author\textsuperscript{1}\thanks{Equal contribution.} \quad
    Second Author\textsuperscript{2}\footnotemark[1] \quad
    Third Author\textsuperscript{1,2}
    \\[0.5em]
    \textsuperscript{1}First Institution, City, Country \\
    \textsuperscript{2}Second Institution, City, Country \\[0.5em]
    \texttt{\{first.author, third.author\}@institution1.edu} \\
    \texttt{second.author@institution2.org}
}

\date{}

% ============================================
% DOCUMENT
% ============================================

\begin{document}

\twocolumn[
\begin{@twocolumnfalse}
\maketitle
\vspace{-1em}
\end{@twocolumnfalse}
]

% ============================================
% ABSTRACT
% ============================================

\begin{abstract}
\noindent Write your abstract here. The abstract should be a single paragraph that summarizes the key contributions and findings of your paper. It should be self-contained and provide readers with a clear understanding of what your paper is about. This template is designed to mimic the style of famous machine learning papers such as DDPM and ``Attention Is All You Need''. It includes commonly used packages, custom commands for mathematical notation, and a clean, professional layout. Keep your abstract between 150--250 words. Avoid citations and undefined acronyms in the abstract.
\end{abstract}


% ============================================
% INTRODUCTION
% ============================================

\section{Introduction}
\label{sec:introduction}

Your introduction goes here. Start with the broader context and motivation for your work. Then narrow down to the specific problem you are addressing.

Key elements to include:
\begin{itemize}[nosep]
    \item Problem statement and motivation
    \item Brief overview of existing approaches and their limitations
    \item Your proposed solution and its advantages
    \item Summary of contributions (often as a bulleted list)
\end{itemize}

\paragraph{Contributions.} Our main contributions are:
\begin{enumerate}[nosep]
    \item First contribution: briefly describe.
    \item Second contribution: briefly describe.
    \item Third contribution: briefly describe.
\end{enumerate}

% ============================================
% RELATED WORK
% ============================================

\section{Related Work}
\label{sec:related_work}

Discuss related work here. Organize by topic or theme rather than listing papers chronologically. Compare and contrast with your approach.

\paragraph{Topic 1.} Discuss relevant work in this area~\citep{example2020paper}.

\paragraph{Topic 2.} Discuss relevant work in this area~\citep{vaswani2017attention}.

% ============================================
% METHOD
% ============================================

\section{Method}
\label{sec:method}

Describe your method here. Start with an overview, then dive into the details.

\subsection{Problem Setup}
\label{sec:problem_setup}

Define your problem formally. For example:

\begin{definition}[Problem Definition]
    Given input data $\x \in \R^d$, we aim to learn a function $f_\theta: \R^d \to \R^k$ that minimizes the expected loss:
    \begin{equation}
        \mathcal{L}(\thetab) = \E_{(\x,\y) \sim p_\text{data}} \left[ \ell(f_{\thetab}(\x), \y) \right]
    \end{equation}
    where $\ell$ is a suitable loss function.
\end{definition}

\subsection{Model Architecture}
\label{sec:architecture}

Describe your model architecture.

\begin{equation}
    \text{Attention}(Q, K, V) = \text{softmax}\left(\frac{QK^\top}{\sqrt{d_k}}\right)V
\end{equation}

\subsection{Training Procedure}
\label{sec:training}

Describe your training procedure. Algorithm pseudocode can be included:

\begin{algorithm}[t]
    \caption{Training Algorithm}
    \label{alg:training}
    \begin{algorithmic}[1]
        \REQUIRE Dataset $\mathcal{D}$, learning rate $\eta$, iterations $T$
        \STATE Initialize parameters $\thetab$
        \FOR{$t = 1$ to $T$}
            \STATE Sample minibatch $\{(\x_i, \y_i)\}_{i=1}^B$ from $\mathcal{D}$
            \STATE Compute loss $\mathcal{L} = \frac{1}{B}\sum_{i=1}^B \ell(f_{\thetab}(\x_i), \y_i)$
            \STATE Update $\thetab \leftarrow \thetab - \eta \nabla_{\thetab} \mathcal{L}$
        \ENDFOR
        \RETURN $\thetab$
    \end{algorithmic}
\end{algorithm}

% ============================================
% EXPERIMENTS
% ============================================

\section{Experiments}
\label{sec:experiments}

\subsection{Experimental Setup}
\label{sec:setup}

Describe your experimental setup, including:
\begin{itemize}[nosep]
    \item Datasets used
    \item Baselines for comparison
    \item Evaluation metrics
    \item Implementation details
\end{itemize}

\subsection{Main Results}
\label{sec:results}

Present your main results. Tables are effective for quantitative comparisons:

\begin{table}[t]
    \centering
    \caption{Results on benchmark datasets. Best in \textbf{bold}.}
    \label{tab:results}
    \small
    \begin{tabular}{lcccc}
        \toprule
        \textbf{Method} & \textbf{A} & \textbf{B} & \textbf{C} & \textbf{Avg.} \\
        \midrule
        Baseline 1 & 85.2 & 78.3 & 91.0 & 84.8 \\
        Baseline 2 & 87.1 & 80.5 & 92.3 & 86.6 \\
        Baseline 3 & 88.4 & 82.1 & 93.1 & 87.9 \\
        \midrule
        Ours & \textbf{91.2} & \textbf{85.7} & \textbf{95.4} & \textbf{90.8} \\
        \bottomrule
    \end{tabular}
\end{table}

\subsection{Ablation Studies}
\label{sec:ablations}

Include ablation studies to validate design choices.

% ============================================
% CONCLUSION
% ============================================

\section{Conclusion}
\label{sec:conclusion}

Summarize your work and contributions. Discuss limitations and future work.

\paragraph{Limitations.} Acknowledge any limitations of your approach.

\paragraph{Future Work.} Suggest directions for future research.

% ============================================
% ACKNOWLEDGMENTS
% ============================================

\section*{Acknowledgments}

We thank the anonymous reviewers for their valuable feedback. This work was supported by [funding sources].

% ============================================
% REFERENCES
% ============================================

\bibliographystyle{plainnat}
\bibliography{references}

% ============================================
% APPENDIX
% ============================================

\appendix

\section{Additional Details}
\label{app:details}

Include supplementary material, proofs, additional experiments, etc.

\section{Proofs}
\label{app:proofs}

\begin{theorem}
    Statement of theorem.
\end{theorem}

\begin{proof}
    Proof goes here.
\end{proof}

\end{document}
